%*****************************************
\chapter{Diskussion}\label{ch:discussion}
%*****************************************
Die Datenerhebung innerhalb der ColoReg Studie, die Grundlage für die Erhebung der hier vorgestellten Ergebnisse bildete, erfolgte unter klinischen Bedingungen.
Daher sind die erhobenen Daten naturgemäß mit Störeinflüssen behaftet, die in einem realen klinischen Umfeld auftreten können.
Diese Störeinflüsse haben sich negativ auf die Qualität der erfassten Daten ausgewirkt und somit die Aussagekraft der Ergebnisse eingeschränkt.
Beispielsweise traten bei einigen Patienten Bewegungsartefakte auf, die die Genauigkeit der Messungen beeinträchtigten.
Zudem gab es Fälle von unvollständigen Datensätzen aufgrund technischer Probleme während der Datenerfassung.
Diese Faktoren führten dazu, dass bei der Evaulierung die vollständige Datenbasis nicht genutzt werden konnte, was die statistische Signifikanz der Ergebnisse reduzierte.
Trotz dieser Einschränkungen konnten jedoch wertvolle Erkenntnisse gewonnen werden, die als Grundlage für zukünftige Studien dienen können.
Zukünftige Untersuchungen sollten darauf abzielen, die Datenerhebung unter kontrollierteren Bedingungen durchzuführen, um die Qualität der Daten zu verbessern und die Aussagekraft der Ergebnisse zu erhöhen.
